\subsection{Confidence Calibration}
\label{subsec:confidence_calibration}

The Confidence Calibration experiment evaluates predictive reliability by investigating the calibration gap between spiking and continuous architectures. Modern \glspl{cnn} are prone to systematic overconfidence, frequently assigning high objectness scores to inaccurate detections \cite{guo2017calibration}. 
Conversely, the temporal dynamics of \glspl{snn} may offer superior uncertainty quantification, providing a more robust reliability profile for safety-critical automotive tasks \cite{skatchkovsky2022bayesian}.

\subsubsection{Methodology: Reliability Diagrams}
To quantify calibration, we generate a Reliability Diagram mapping predicted confidence ($C$) against actual regression accuracy ($IoU$) for every validation detection.
In a perfectly calibrated system, these values align along the identity line ($C = IoU$). 
Clumping in the high-confidence, low-accuracy quadrant indicates an overconfidence bias, where the model fails to signal its own predictive errors.

\subsubsection{Hypothesis: Spikes as Uncertainty Sensors}
The \gls{cnn} baseline is hypothesised to exhibit significant miscalibration due to its continuous activation functions and exact gradient descent, which tend to overfit the confidence head. 
The \gls{snn} is expected to demonstrate a more linear, calibrated trend. 
In the spiking domain, low-quality features generate fewer spikes, leading to lower membrane potentials in the detection head. 
This mechanism allows the \gls{snn} to function as an implicit uncertainty sensor, where discrete spike rates effectively gate the output confidence logit relative to the strength of the temporal evidence.
